\documentclass[11pt]{article}

\usepackage[margin=1in]{geometry}
\usepackage{amsmath,amssymb,amsthm}
\usepackage{booktabs}
\usepackage{enumitem}
\usepackage[hidelinks]{hyperref}

\newtheorem{theorem}{Theorem}
\newtheorem{proposition}[theorem]{Proposition}
\newtheorem{corollary}[theorem]{Corollary}
\newtheorem{lemma}[theorem]{Lemma}
\theoremstyle{definition}
\newtheorem{definition}[theorem]{Definition}
\theoremstyle{remark}
\newtheorem{remark}[theorem]{Remark}

\newcommand{\Z}{\mathbb Z}
\newcommand{\F}{\mathbb F}
\newcommand{\ShUg}{\operatorname{ShUg}}
\newcommand{\Arf}{\operatorname{Arf}}
\newcommand{\taglabel}[1]{\textnormal{\textbf{[#1]}}}
\newcommand{\PROVED}{\taglabel{PROVED}}
\newcommand{\SOURCE}{\taglabel{SOURCE-PINNED}}
\newcommand{\BOUNDARY}{\taglabel{BOUNDARY}}

\title{Brown phases in game semantics:\\a canonical two-bit split and a divisibility obstruction}
\author{Codex}
\date{2026-08-09}

\begin{document}
\maketitle

\begin{abstract}
The Brown invariant of a nonsingular $\Z/4$-valued quadratic form is the
eighth-root phase of a four-class Gauss census.  We determine exactly what a
game-theoretic reading of that census adds to the characteristic-two Arf
problem.  Every Brown form has a canonical decomposition
\[
 q=\widehat\ell+2Q,
\]
where $\ell=q\bmod2$ is linear and $Q$ is an ordinary $\F_2$-quadratic
form.  Consequently the four residue classes are precisely the joint census
of one linear game bit and the ordinary quadratic bit already isolated by the
\texttt{tis} problem; the Brown phase introduces no second nonlinear
game-semantic problem.  We also prove a complementary global no-go: an
ambient-coherent Brown refinement of the additive group of all short games is
identically zero, because that group is $2$-divisible.  Finally, the cyclic
extension $C_4\to C_8\to C_2$ shows that a bare central extension by $C_4$
does not determine a Brown form: a section or phase framing is essential.
Together these results dispose of the independent Brown invariant/output
question under an explicit synchronized-charge contract, without claiming a
single-game outcome construction or a solution of \texttt{tis} itself.
\end{abstract}

\section{The exact question}

A Brown form on an $\F_2$-space $V$ is a function
\[
 q:V\longrightarrow\Z/4
\]
with $q(0)=0$ and
\begin{equation}\label{eq:brown-law}
 q(x+y)=q(x)+q(y)+2b(x,y),
\end{equation}
where $b:V\times V\to\F_2$ is symmetric and bilinear.  Substituting
$y=x$ shows
\begin{equation}\label{eq:diag}
 b(x,x)=q(x)\bmod2.
\end{equation}
Thus $b$ need not be alternating.  When $b$ is nonsingular, the classical
Brown invariant $\beta(q)\in\Z/8$ is defined by the Gauss sum
\begin{equation}\label{eq:gauss}
 \sum_{x\in V} i^{q(x)}=2^{\dim(V)/2}\zeta_8^{\beta(q)}.
\end{equation}
This is standard mathematics due to Brown and Wall
\cite{brown,wall}; the repository's \texttt{brown\_f2} implementation also
records the singular radical boundary.

The phrase ``game-theoretic reading'' has three logically different senses:
\begin{enumerate}[label=(\roman*),leftmargin=2.5em]
\item a four-class charge or outcome attached to positions in one selected
      finite $\F_2$-space;
\item a quadratic colour attached intrinsically to short-game \emph{values}
      and compatible with disjunctive sum and ambient enlargement; or
\item a framed evaluator which is handed $q$ and reports its residue.
\end{enumerate}
Without separating these, four raw labels prove nothing: partizan normal play,
normal/misere pairs, loopy starter pairs, and scoring games all provide at
least four labels.  The results below reduce (i) algebraically to a
synchronized linear/quadratic charge pair, rule out (ii) nontrivially, and
identify (iii) as tautological rather than a new bridge.

\section{The canonical two-bit decomposition}

Write $\rho:\Z/4\to\F_2$ for reduction modulo two and
$\widehat{\phantom a}:\F_2\to\Z/4$ for the distinguished set-theoretic
section $0\mapsto0$, $1\mapsto1$.  It obeys
\begin{equation}\label{eq:carry}
 \widehat{a+c}=\widehat a+\widehat c+2ac
 \qquad(a,c\in\F_2),
\end{equation}
where the left addition is in $\F_2$ and the right addition is in $\Z/4$.

\begin{theorem}[canonical Brown split]\label{thm:split}
\PROVED{} Let $(q,b)$ satisfy \eqref{eq:brown-law}.  There are unique maps
\[
 \ell:V\to\F_2,
 \qquad
 Q:V\to\F_2
\]
such that
\begin{equation}\label{eq:split}
 q(x)=\widehat{\ell(x)}+2Q(x).
\end{equation}
They have the following properties.
\begin{enumerate}[label=(\roman*),leftmargin=2.5em]
\item $\ell=\rho\circ q$ is linear.
\item $Q$ is an ordinary characteristic-two quadratic form with alternating
      polar form
      \begin{equation}\label{eq:new-polar}
       B_Q(x,y)=b(x,y)+\ell(x)\ell(y).
      \end{equation}
\item Conversely, every linear $\ell$ and every ordinary $\F_2$-quadratic
      form $Q$ determine a unique Brown form by \eqref{eq:split}; its symmetric
      polar form is
      \[
       b=B_Q+\ell\otimes\ell.
      \]
Hence Brown forms are canonically equivalent to pairs $(\ell,Q)$.
\end{enumerate}
\end{theorem}

\begin{proof}
Reducing \eqref{eq:brown-law} modulo two kills $2b$, so
$\ell=\rho q$ is additive and therefore $\F_2$-linear.  Each residue
$q(x)$ has a unique expression \eqref{eq:split}, which defines $Q$ and proves
uniqueness.

Apply \eqref{eq:carry} to $\ell(x+y)=\ell(x)+\ell(y)$ and compare the two
expressions for $q(x+y)$.  Cancelling the distinguished lifts and dividing
the remaining even equality by two gives
\[
 Q(x+y)=Q(x)+Q(y)+b(x,y)+\ell(x)\ell(y).
\]
Thus the displayed $B_Q$ is the polar form of $Q$.  Equation
\eqref{eq:diag} and $\ell(x)^2=\ell(x)$ give
\[
 B_Q(x,x)=b(x,x)+\ell(x)^2=\ell(x)+\ell(x)=0,
\]
so it is alternating; bilinearity and symmetry follow from those of $b$ and
$\ell\otimes\ell$.

Conversely, let $\ell$ be linear and $Q$ have alternating polar $B_Q$.
Set $q=\widehat\ell+2Q$ and $b=B_Q+\ell\otimes\ell$.  Equation
\eqref{eq:carry} and the polarization law for $Q$ give
\eqref{eq:brown-law}; moreover $b(x,x)=\ell(x)=q(x)\bmod2$.
The two constructions are inverse by the uniqueness of the two binary digits
of every element of $\Z/4$.
\end{proof}

\begin{corollary}[the four-class census is a joint binary census]
\label{cor:census}
\PROVED{} Under Theorem~\ref{thm:split}, the four Brown classes are
\begin{center}
\begin{tabular}{cccc}
\toprule
$q(x)$ & $\ell(x)$ & $Q(x)$ & $i^{q(x)}$\\
\midrule
$0$ & $0$ & $0$ & $+1$\\
$1$ & $1$ & $0$ & $+i$\\
$2$ & $0$ & $1$ & $-1$\\
$3$ & $1$ & $1$ & $-i$\\
\bottomrule
\end{tabular}
\end{center}
In particular
\begin{equation}\label{eq:joint-gauss}
 \sum_x i^{q(x)}=\sum_x i^{\ell(x)}(-1)^{Q(x)}.
\end{equation}
If $W(R)=\sum_x(-1)^{R(x)}$, then more explicitly
\begin{equation}\label{eq:walsh-pair}
 \sum_x i^{q(x)}
 =\frac{1+i}{2}W(Q)+\frac{1-i}{2}W(Q+\ell).
\end{equation}
When the full Gauss sum is nonzero---in particular when $b$ is
nonsingular---$\beta$ is the octant of the \emph{joint} $(\ell,Q)$
imbalance, or equivalently of the correlated pair of ordinary quadratic
signed censuses $Q$ and $Q+\ell$.  It is not generally the product of a
marginal linear census and the Arf bias of $Q$.  On the doubled nonsingular
slice $\ell=0$, the two Walsh sums agree and this reduces to the shipped
identity $\beta(2Q)=4\Arf(Q)$.  For singular forms the same census identity
holds, while the phase belongs to the nonsingular core and the radical data
is recorded separately.
\end{corollary}

\begin{proof}
The table and \eqref{eq:joint-gauss} follow immediately from
$q=\widehat\ell+2Q$.  The identity
\[
 i^{\ell}=\frac{1+i}{2}+\frac{1-i}{2}(-1)^{\ell}
\]
gives \eqref{eq:walsh-pair}.  For $\ell=0$, the Gauss sum is the ordinary
real quadratic Gauss sum, whose sign is $(-1)^{\Arf(Q)}$.
\end{proof}

\subsection{The Brown phase in ordinary quadratic coordinates}

The split gives more than a four-colour relabelling: it computes every Brown
octant from ordinary characteristic-two data.

\begin{theorem}[phase formulas]\label{thm:phase}
\PROVED{} Suppose the symmetric Brown polar form $b$ is nondegenerate and put
$n=\dim V$.
\begin{enumerate}[label=(\roman*),leftmargin=2.5em]
\item If $n$ is even, $B_Q$ is nondegenerate.  Let $a$ be the unique vector
      satisfying $\ell(x)=B_Q(a,x)$ for every $x$.  Then
      \begin{equation}\label{eq:even-beta}
       \beta(q)=4\Arf(Q)+2Q(a)\pmod8.
      \end{equation}
\item If $n$ is odd, $\operatorname{rad}B_Q=\langle w\rangle$ is
      one-dimensional, $\ell(w)=1$, and $B_Q$ restricts nondegenerately to
      $H=\ker\ell$.  Then
      \begin{equation}\label{eq:odd-beta}
       \beta(q)=4\Arf(Q|_H)+
       \begin{cases}
       1,&Q(w)=0,\\
       7,&Q(w)=1
       \end{cases}
       \pmod8.
      \end{equation}
\end{enumerate}
\end{theorem}

\begin{proof}
For a nonsingular symmetric form over $\F_2$, the characteristic functional
$x\mapsto b(x,x)$ has a unique representing vector $c$.  The standard
orthogonal line/hyperbolic-plane decomposition shows
$b(c,c)=n\bmod2$.  Since $\ell(x)=b(c,x)$ and
$B_Q=b+\ell\otimes\ell$, its kernel equation is
\[
 B_Q(v,-)=0
 \quad\Longleftrightarrow\quad
 v=\ell(v)c.
\]
Thus $B_Q$ is nonsingular when $n$ is even; when $n$ is odd its radical is
$\langle c\rangle$, with $\ell(c)=1$, and its restriction to
$H=\ker\ell$ is nonsingular.  Put $w=c$ in the odd case.

In even dimension, nondegeneracy supplies the unique $a$ with
$B_Q(a,-)=\ell$.  Translation by $a$ gives
\[
 Q(x+a)=Q(x)+Q(a)+\ell(x),
\]
and hence $W(Q+\ell)=(-1)^{Q(a)}W(Q)$.  The ordinary quadratic Gauss
formula gives $W(Q)=(-1)^{\Arf(Q)}2^{n/2}$.  Substitution in
\eqref{eq:walsh-pair} yields \eqref{eq:even-beta}.

In odd dimension, pair the two cosets as $h$ and $h+w$, $h\in H$.  Since
$w\in\operatorname{rad}B_Q$,
\[
 \sum_{x\in V}i^{q(x)}
 =\sum_{h\in H}(-1)^{Q(h)}
   \bigl(1+i(-1)^{Q(w)}\bigr).
\]
The first factor is
$(-1)^{\Arf(Q|_H)}2^{(n-1)/2}$; the second has phase $\zeta_8$ or
$\zeta_8^7$.  This proves \eqref{eq:odd-beta}.
\end{proof}

\begin{remark}[degenerate forms]
Let $R=\operatorname{rad}b$.  Equation~\eqref{eq:diag} gives
$\ell|_R=0$, hence $q|_R=2Q|_R$.  If $Q|_R\ne0$, translation by a vector
with $Q=1$ cancels the full Gauss sum.  If $Q|_R=0$, the pair $(\ell,Q)$
descends to $V/R$ and Theorem~\ref{thm:phase} applies to the nonsingular
quotient.  This is exactly the repository's
\texttt{radical\_anisotropic} boundary.
\end{remark}

\begin{remark}[the split is twisted under addition]
Theorem~\ref{thm:split} is a canonical bijection of data, not a direct-product
law.  Under the Brown--Wall orthogonal sum, if $q_1,q_2$ have coordinates
$(\ell_1,Q_1)$ on $V_1$ and $(\ell_2,Q_2)$ on $V_2$, then on
$V_1\oplus V_2$
\begin{equation}\label{eq:carry-law}
 \begin{aligned}
 \ell(x_1,x_2)&=\ell_1(x_1)+\ell_2(x_2),\\
 Q(x_1,x_2)&=Q_1(x_1)+Q_2(x_2)
   +\ell_1(x_1)\ell_2(x_2).
 \end{aligned}
\end{equation}
The same formula holds for pointwise addition on one source whenever that
sum is considered.  The last term is the binary carry from
\eqref{eq:carry}: it prevents this coordinate split from being a
direct-product law and reflects the cited Brown--Wall $\Z/8$ Witt structure;
it is not by itself a proof of that classification.  A positionwise joint
readout retains the carry and the correlation; separate marginal counts do
not.
\end{remark}

\section{What the split means for play}

\begin{definition}[charge readout]
A binary or four-class \emph{charge readout} is a game rule whose complete
play assigns a terminal element of $\F_2$ or $\Z/4$.  It is weaker than a
normal-play $P/N$ realization: the terminal charge may be strategically forced
without itself being the winner class.
\end{definition}

\begin{definition}[synchronized four-charge contract]
Two charge readouts on the same input may be observed jointly and their
ordered pair may be relabelled by a fixed bijection.  Equivalence in this
contract means equality of the resulting terminal labels for every input and
complete play.  It does \emph{not} assert that the two game trees factor, that
their strategies are independent, or that the pair is already one of the
canonical outcome classes of a single normal-, misere-, or loopy-play game.
\end{definition}

\begin{proposition}[no independent \texttt{over} charge]
\label{prop:game-split}
\PROVED{} A four-class charge readout realizing a Brown form $q$ is equivalent,
by deterministic relabelling, to the joint readout $(\ell,Q)$ of
Theorem~\ref{thm:split}.  Conversely, joint readouts of $\ell$ and $Q$ give a
four-class readout by $(a,c)\mapsto\widehat a+2c$, provided the semantics
admits synchronized products of two charge channels on the same input.

Once $\ell$ is supplied, its coordinate is a solved local linear charge
layer: after choosing any basis, attach a $*$ component to precisely those
active coordinates on which $\ell$ is one; their disjunctive sum has Grundy
value $*\ell(x)$.  All nonlinear dependence of the terminal charge on the
input is therefore the ordinary quadratic readout $Q$.
If one strengthens ``charge readout'' to ``$P$ exactly when $Q=0$'', that
remaining requirement is a generalized \texttt{tis}-type problem; it is
literally the existing Gold-form \texttt{tis} problem on the doubled slice
$\ell=0$, $q=2Q$.
\end{proposition}

\begin{proof}
The residue-to-bit map is the table in Corollary~\ref{cor:census}, and its
inverse is $(a,c)\mapsto\widehat a+2c$; this is an observational equivalence
of terminal labels, not a factorization theorem for game trees.
For the linear realization, write $x=\sum_i x_i e_i$.  The sum of the tagged
$*$ components has Grundy coordinate
$\sum_i x_i\ell(e_i)=\ell(x)$.  Changing basis changes the presentation, not
the resulting linear functional.  The final sentence is the definition of
the \texttt{tis} target applied to the ordinary quadratic form $Q$.
\end{proof}

\begin{remark}[the exact residual boundary]
\BOUNDARY{} Proposition~\ref{prop:game-split} does not turn a forced terminal
charge into a normal-play winner.  The repository's $\sigma$-valued
echo--FIFO construction is therefore genuine progress toward the $Q$ bit but
does not, without the general linking theorem and an outcome recasting, solve
\texttt{tis}.  Nor is a synchronized ordered pair automatically one canonical
normal-, misere-, or loopy-play outcome of a single game; product closure is
part of the four-charge contract.  The point is narrower and complete:
Brown's four classes add a linear axis and a correlated second census, not a
new algebraic kind of nonlinear charge/output problem.
\end{remark}

\section{The global value-level obstruction}

The previous section concerns a selected finite $\F_2$-space.  A stronger
naturality demand is that $q$ belong to short-game values themselves rather
than to one root-incomplete slice.

\begin{definition}[ambient-coherent Brown value datum]
For every finitely generated subgroup $M\leq\ShUg$, choose
\[
 q_M:M\to\Z/4,
 \qquad b_M:M\times M\to\F_2
\]
satisfying \eqref{eq:brown-law}, with $b_M$ biadditive.  The family is
\emph{ambient-coherent} if restriction along every inclusion $M\subseteq N$
preserves both functions.  This is a generalized Brown-law datum on arbitrary
game groups, deliberately stronger than a classical Brown form whose domain
is elementary abelian.
\end{definition}

\begin{theorem}[short-game Brown collapse]\label{thm:global}
\PROVED{} Every ambient-coherent Brown value datum on short games is zero:
\[
 b_M=0,
 \qquad q_M=0
\]
for every finitely generated $M\leq\ShUg$.
\end{theorem}

\begin{proof}
Moews proves
\begin{equation}\label{eq:moews}
 \ShUg\cong\bigoplus_{\mathbb N}\Z[1/2]
 \oplus\bigoplus_{\mathbb N}\bigl(\Z[1/2]/\Z\bigr),
\end{equation}
so multiplication by two, and hence by four, is surjective
\cite[Theorem~9]{moews}.

Fix $x,y\in M$.  Choose $u\in\ShUg$ with $2u=x$ and enlarge $M$ to the
finitely generated subgroup $N=\langle M,u\rangle$.  Coherence and
bilinearity give
\[
 b_M(x,y)=b_N(2u,y)=2b_N(u,y)=0\quad\text{in }\F_2.
\]
Thus all polar corrections vanish and every $q_M$ is additive.  Choose
$v\in\ShUg$ with $4v=x$ and put $P=\langle M,v\rangle$.  Then
\[
 q_M(x)=q_P(x)=q_P(4v)=4q_P(v)=0\quad\text{in }\Z/4.
\]
\end{proof}

\begin{corollary}[no exponent-two Brown quotient]
\PROVED{} Every additive map from $\ShUg$ to an exponent-two group is zero;
in particular $\ShUg/2\ShUg=0$.  Hence no nonzero $\F_2$ Brown module can
be an additive quotient or homomorphic image of the full short-game group.
This does not kill the intrinsic subgroup $\ShUg[2]$: it contains $*$ and
its root-incomplete finite subspaces support nonzero local Brown forms.
\end{corollary}

\begin{proof}
If $f:\ShUg\to A$ is additive, $2A=0$, and $x=2u$, then
$f(x)=2f(u)=0$.
\end{proof}

This theorem is stronger than the earlier game-exterior obstruction on this
specific target.  It needs no Clifford algebra, coefficient embedding, or
torsion hypothesis: the Brown polarization law and ambient roots suffice.
The kernel-checked file \texttt{formal/Ogdoad/BrownGame.lean} proves the
abstract two-divisible/exponent-four implication and the $\Z/4$
specialization; Equation~\eqref{eq:moews} remains an explicit source-pinned
input rather than a Lean axiom.

\section{Why a bare central extension does not supply the phase}

The characteristic-two extraspecial construction does not lift verbatim from
$C_2$ to $C_4$.  The obstruction is not commutativity but missing section
data.

\begin{proposition}[minimal framed cyclic model]\label{prop:c8}
\PROVED{} Consider the abelian exact sequence
\[
 0\longrightarrow C_4\xrightarrow{\,a\mapsto2a\,}C_8
 \longrightarrow C_2\longrightarrow0.
\]
Choosing $1\in C_8$ as the lift of the nonzero quotient element gives the odd
Brown line $q(1)=1$; choosing $3\in C_8$ gives $q(1)=3$.  Both have
$b(1,1)=1$ and Brown phases $1$ and $7$, respectively.  Therefore the bare
extension determines neither $q$ nor $\beta$.
\end{proposition}

\begin{proof}
Let $z=2\in C_8$ generate the embedded $C_4$.  For the first section,
$2\cdot1=z^1$ in additive notation, while for the second,
$2\cdot3=6=z^3$.  In either case, substituting the nonzero vector twice in
\eqref{eq:brown-law} forces $b(1,1)=1$.  The Gauss sums are $1+i$ and
$1-i$, hence the stated phases.  More generally, changing a lift $g$ to
$z+g$ changes its double by $2z$; the square exponent is not fibre-invariant.
\end{proof}

The finite model and the non-fibre-invariance statement are also
kernel-checked in \texttt{BrownGame.lean}.  It corrects the tempting but false
claim that a $C_4$-central extension ``corresponds exactly'' to a Brown form.
What corresponds to the full refinement is a suitably section-framed
extension.  Its commutator remembers only alternating information; the
admissible section/framing carries the missing quadratic lift data, including
the diagonal phase.

\section{Disposition of \texttt{over}}

The three readings from Section~1 now have exact answers.
\begin{enumerate}[label=(\arabic*),leftmargin=2.5em]
\item \PROVED{} On a selected finite $\F_2$-space, the Brown four-class
      census is canonically the joint census of a locally realizable linear
      charge $\ell$ (once supplied) and one ordinary quadratic bit $Q$.  Its complex imbalance is
      \eqref{eq:joint-gauss}.
\item \PROVED{} On short-game values with ambient coherence, every such
      Brown datum is zero by Theorem~\ref{thm:global}.
\item \PROVED{} A bare $C_4$-central extension is insufficient; a chosen
      section can encode the phase, but then the framing is supplied data.
\end{enumerate}

Accordingly, under the explicit synchronized four-charge contract, there is
no independent mod-$8$ charge/output problem left.  A non-tautological
autonomous rule for the nonlinear bit is a generalized \texttt{tis}-type
problem; on doubled Gold forms it is literally the surviving \texttt{tis}
problem.  Crossing such a readout with the linear bit gives the Brown classes
automatically.  This closes \texttt{over} as a separate invariant/output
direction on that contract without promoting the verified $\sigma$-charge or
finite census to a proof of \texttt{tis}, and without asserting that the
synchronized pair is already a single canonical normal/misere/loopy outcome
semantics.  If that single-game internalization is demanded, it is the exact
residual boundary and has been transferred rather than proved away.

\begin{thebibliography}{9}
\bibitem{brown}
Edgar H. Brown, Jr.,
\emph{Generalizations of the Kervaire invariant},
Annals of Mathematics \textbf{95} (1972), 368--383.
\url{https://doi.org/10.2307/1970804}.

\bibitem{moews}
David Moews,
\emph{The Abstract Structure of the Group of Games},
in \emph{More Games of No Chance}, MSRI Publications 42 (2002), 49--58.
\url{https://library.slmath.org/books/Book42/files/moews.pdf}.

\bibitem{wall}
C.~T.~C. Wall,
\emph{Quadratic forms on finite groups, and related topics},
Topology \textbf{2} (1963), 281--298.
\url{https://doi.org/10.1016/0040-9383(63)90012-0}.
\end{thebibliography}

\end{document}
